\DocumentMetadata{
  pdfversion=2.0,
  pdfstandard=ua-2,
  lang=en-US,
  tagging=on,
  tagging-setup={math/setup=mathml-SE,tabsorder=structure}
}
\documentclass[11pt,letterpaper]{article}
\usepackage[margin=0.68in,includefoot]{geometry}
\usepackage{newcomputermodern}
\usepackage{amsmath,amscd,mathtools}
\providecommand{\square}{\mdlgwhtsquare}
\usepackage[hidelinks]{hyperref}
\hypersetup{
  pdftitle={Algebra PhD Exam, August 2005},
  pdfauthor={University of Florida Department of Mathematics},
  pdfsubject={Graduate mathematics examination},
  pdfdisplaydoctitle=true,
  bookmarksopen=true
}
\setcounter{secnumdepth}{0}
\setlength{\parindent}{0pt}
\setlength{\parskip}{0.28em}
\setlength{\emergencystretch}{3em}
\setlength{\leftmargini}{2.15em}
\setlength{\leftmarginii}{2.65em}
\setlength{\leftmarginiii}{3em}
\renewcommand{\labelenumi}{\textbf{\theenumi.}}
\pagestyle{empty}
\raggedbottom
\allowdisplaybreaks
\newenvironment{examproblems}[1]{%
  \begin{enumerate}%
  \setcounter{enumi}{#1}%
  \setlength{\topsep}{0.35em}%
  \setlength{\partopsep}{0pt}%
  \setlength{\itemsep}{0.48em}%
  \setlength{\parsep}{0.18em}%
}{\end{enumerate}}
\newenvironment{examsubparts}{%
  \begin{enumerate}%
  \setlength{\topsep}{0.18em}%
  \setlength{\partopsep}{0pt}%
  \setlength{\itemsep}{0.16em}%
  \setlength{\parsep}{0.1em}%
}{\end{enumerate}}
\newenvironment{examinstructions}{%
  \par\begingroup\itshape
}{%
  \par\endgroup\vspace{0.18em}
}
\makeatletter
\renewcommand\section{\@startsection{section}{1}{0pt}%
  {-1.25ex plus -.25ex minus -.1ex}%
  {0.55ex plus .1ex}%
  {\normalfont\large\bfseries}}
\renewcommand{\@maketitle}{%
  \newpage\null\vskip -1.2em%
  {\footnotesize\bfseries
    \href{https://www.ufl.edu/}{University of Florida}, \href{https://math.ufl.edu/}{Department of Mathematics}\par}%
  \vskip 0.18em%
  \tagstructbegin{tag=Title}%
    {\LARGE\bfseries\href{https://gma.math.ufl.edu/exams/algebra-phd/}{Algebra PhD Exam}, August 2005\par}%
  \tagstructend%
  \vskip 0.35em%
  \tagmcbegin{artifact}%
    \hrule height 1pt%
  \tagmcend%
  \vskip 0.55em%
}
\makeatother
\title{Algebra PhD Exam, August 2005}
\author{}
\date{}
\begin{document}
\maketitle
\thispagestyle{empty}
\begin{examinstructions}
Answer seven problems. Write your answers clearly in complete English sentences. You may quote results (within reason) as long as you state them clearly.
\end{examinstructions}
\begin{examproblems}{0}
\item
Let \(\mathcal C\) be a category and let~\(\Lambda\) be a set. For each \(\lambda\in\Lambda\), let \(X_\lambda\) be an object in \(\mathcal C\).
\begin{examsubparts}
\item[\textnormal{(a)}]
Give the definition of a product of the collection of objects \(\{X_\lambda\}_{\lambda\in\Lambda}\).
\item[\textnormal{(b)}]
Prove that any two products of the \(\{X_\lambda\}_{\lambda\in\Lambda}\) are isomorphic.
\item[\textnormal{(c)}]
Prove that products always exist in the category of groups.
\end{examsubparts}
\item
Show that the alternating group \(A_n\) is a simple group for \(n\geq 5\). Show that the symmetric group \(S_n\) is not solvable for \(n\geq 5\).
\item
State and prove Jacobson’s Density Theorem.
\item
Consider the polynomial 
\[
p(x)=x^5-3x^4-2x^2+2x+2
\]
 over \(\mathbb Q\). Can all the roots of \(p(x)\) be written by radicals over \(\mathbb Q\)? Can they be written by radicals over \(\mathbb C\)? Carefully justify both answers.
\item
Suppose that~\(F\) is a finite extension of \(\mathbb Q\). Prove that~\(F\) contains only a finite number of roots of unity.
\item
Let~\(F\) be the free group in generators~\(x\) and~\(y\). Let~\(N\) be the subgroup of~\(F\) generated by the set 
\[
S=\{z^2:z\in F\}.
\]
 Prove that~\(N\) is a normal subgroup of~\(F\), and calculate the order of \(F/N\).
\item
Let~\(A\) be a torsion abelian group, and let \(\mathbb Q\) be the (additive) group of the rational numbers. Prove that \(A\otimes_{\mathbb Z}\mathbb Q=\{0\}\), and that \(\mathbb Q\otimes_{\mathbb Z}\mathbb Q\cong\mathbb Q\).
\item
Let~\(R\) be an integral domain with quotient field~\(F\). Suppose that \(a\in R\) is such that \(a\neq 0\) and \(1/a\) is integral over~\(R\). Prove that~\(a\) is a unit of~\(R\).
\item
Let~\(R\) be a commutative ring with \(1\neq 0\). Prove that the set consisting of~\(0\) and all zero divisors in~\(R\) contains at least one prime ideal.
\item
Let~\(R\) be a commutative ring with identity different from zero. Let \(F_1\) and \(F_2\) be two finitely generated free~\(R\)-modules, and assume that \(F_1\cong F_2\) as~\(R\)-modules. Prove that the cardinality of the set of free generators for \(F_1\) equals the cardinality of the set of free generators for \(F_2\).
\item
Let~\(R\) and~\(S\) be semi-simple rings with exactly \(400\) elements. Assume that neither~\(R\) nor~\(S\) is commutative. Does it follow that \(R\cong S\)? Justify your answer.
\end{examproblems}
\end{document}
